% Figure 1: functional roles of generative models in an HRI pipeline
\begin{figure}[t]
\centering
\resizebox{\textwidth}{!}{%
\begin{tikzpicture}[
    font=\sffamily\footnotesize,
    role/.style={draw=black!70, rounded corners=2pt, minimum height=11mm,
                 text width=23mm, align=center, fill=black!4, inner sep=2pt},
    inst/.style={draw=black!70, dashed, rounded corners=2pt, minimum height=9mm,
                 text width=27mm, align=center, fill=white, inner sep=2pt},
    actor/.style={draw=black!70, thick, minimum height=11mm, text width=15mm,
                  align=center, fill=white},
    flow/.style={-{Latex[length=2mm]}, semithick, black!75},
    obs/.style={-{Latex[length=2mm]}, semithick, dashed, black!60},
    band/.style={font=\sffamily\scriptsize\itshape, text=black!70}]

% --- actors --------------------------------------------------------------
\node[actor] (human) at (0,0) {Person(s)};
\node[actor] (human2) at (14.6,0) {Person(s)};

% --- in-robot roles ------------------------------------------------------
\node[role] (perc) at (2.65,0) {\textbf{R1 Perceiver}\\interprets people\\and scenes};
\node[role] (rep)  at (5.45,0) {\textbf{R2 Representer}\\embeds, retrieves,\\remembers};
\node[role] (plan) at (8.25,0) {\textbf{R3 Decider}\\plans, selects,\\times, defers};
\node[role] (lang) at (11.4,1.45) {\textbf{R4 Speaker}\\user-facing\\language};
\node[role] (beh)  at (11.4,0) {\textbf{R5 Behaviour author}\\expressive behaviour,\\code, parameters};
\node[role] (act)  at (11.4,-1.45) {\textbf{R6 Actor}\\end-to-end\\physical action};

\draw[flow] (human) -- (perc);
\draw[flow] (perc) -- (rep);
\draw[flow] (rep) -- (plan);
\draw[flow] (plan.east) -- ++(0.35,0) |- (lang.west);
\draw[flow] (plan.east) -- (beh.west);
\draw[flow] (plan.east) -- ++(0.35,0) |- (act.west);
\draw[flow] (lang.east) -- ++(0.45,0) |- ([yshift=2.5mm]human2.west);
\draw[flow] (beh.east) -- (human2.west);
\draw[flow] (act.east) -- ++(0.45,0) |- ([yshift=-2.5mm]human2.west);

% robot boundary
\begin{scope}[on background layer]
  \node[draw=black!55, thick, rounded corners=4pt, fill=black!1,
        fit=(perc)(rep)(plan)(lang)(beh)(act), inner sep=3.2mm] (robot) {};
\end{scope}
\node[band, anchor=south west] at ([xshift=1mm,yshift=0.2mm]robot.north west)
     {Robot architecture (models deployed in the interaction loop)};

% end-to-end bypass (VLA)
\draw[flow, black!45] (perc.south) |- (act.west |- 0,-2.05) -| ([xshift=-6mm]act.south);
\node[band] at (6.1,-1.86) {end-to-end vision--language--action policies collapse R1--R3 into R6};

% --- research-instrument roles ------------------------------------------
\node[inst] (anno) at (3.3,-3.75) {\textbf{R7 Annotator}\\codes transcripts,\\video, logs};
\node[inst] (judge) at (7.3,-3.75) {\textbf{R8 Judge}\\scores outputs of\\R1--R6};
\node[inst] (sim) at (11.3,-3.75) {\textbf{R9 Simulator}\\synthetic users,\\partners, raters};

\begin{scope}[on background layer]
  \node[draw=black!55, thick, dashed, rounded corners=4pt,
        fit=(anno)(judge)(sim), inner sep=3.2mm] (lab) {};
\end{scope}
\node[band, anchor=north west] at ([xshift=1mm,yshift=-0.2mm]lab.south west)
     {Research pipeline (models used to study the interaction)};

\draw[obs] (anno.north) -- (anno.north |- robot.south);
\draw[obs] (judge.north) -- (judge.north |- robot.south);
\draw[obs] (sim.north) -- (sim.north |- robot.south);
\end{tikzpicture}}
\caption{Nine functional roles that language, vision, and action models play in
HRI research. Solid boxes (R1--R6) are roles inside the robot's interaction
loop; dashed boxes (R7--R9) are roles inside the research pipeline that
observes, scores, or stands in for the interaction. A single model usually
occupies several roles at once, and each role inherits a different evaluation
tradition (Table~\ref{tab:roles}).}
\Description{A block diagram. A person on the left feeds into a chain of three
boxes inside a robot boundary labelled R1 Perceiver, R2 Representer and R3
Decider. The Decider fans out to three output boxes, R4 Speaker, R5 Behaviour
author and R6 Actor, each of which points to a person on the right. A grey
arrow bypasses the chain from the Perceiver to the Actor to indicate
end-to-end vision-language-action policies. Below the robot boundary, a dashed
boundary labelled research pipeline contains three dashed boxes, R7 Annotator,
R8 Judge and R9 Simulator, with dashed arrows pointing up at the robot.}
\label{fig:roles}
\end{figure}
